% =============================================================================
%  Reinforcement Learning for Robotics — The FCP Way
%  Print edition, in the Shannon Robotics design language.
%
%  Build:  latexmk -pdf main.tex      (or: pdflatex main.tex, three times)
% =============================================================================
\documentclass[11pt, twoside, openright]{book}

\usepackage[
  paperwidth=7in, paperheight=10in,
  inner=0.95in, outer=0.75in, top=0.85in, bottom=0.95in,
  headsep=16pt, footskip=28pt
]{geometry}

\usepackage{shannon}

\hypersetup{
  pdftitle={Reinforcement Learning for Robotics --- The FCP Way},
  pdfauthor={Shannon Robotics},
  pdfsubject={Reinforcement learning, robotics, Rust},
  colorlinks=true,
  linkcolor=shAccent,
  urlcolor=shSkyDark,
  citecolor=shAccent
}

\begin{document}

% =============================================================================
%  Front matter
% =============================================================================
\frontmatter
\pagestyle{empty}

\begin{titlepage}
\thispagestyle{empty}
\noindent
{\color{shTeal}\rule{\textwidth}{3pt}}
\vspace{2.2in}

\noindent
{\displayfont\fontsize{10}{12}\selectfont\color{shTealDark}%
 \MakeUppercase{Shannon Robotics}}

\vspace{10pt}
\noindent
{\displayfont\bfseries\fontsize{34}{39}\selectfont\color{shInk}%
 Reinforcement Learning\\for Robotics}

\vspace{12pt}
\noindent
{\displayfont\fontsize{16}{20}\selectfont\color{shInkMuted}%
 The FCP Way}

\vspace{20pt}
\noindent
{\color{shSlate300}\rule{2.4in}{0.6pt}}

\vspace{16pt}
\noindent
\begin{minipage}{4.2in}
{\fontsize{11}{15}\selectfont\color{shInkSecondary}%
 Complete mathematical foundations, interactive intuition, and working
 implementations in Rust --- for robots that must act in a world that does not
 hold still.}
\end{minipage}

\vfill
\noindent
{\displayfont\fontsize{9}{12}\selectfont\color{shInkMuted}%
 Foundation \ \textperiodcentered\ \ Conceptual \ \textperiodcentered\ \ Practical}

\vspace{6pt}
\noindent
{\color{shTeal}\rule{\textwidth}{3pt}}
\end{titlepage}

% ----- Colophon --------------------------------------------------------------
\cleardoublepage
\thispagestyle{empty}
\vspace*{\fill}
\noindent
{\fontsize{9}{13}\selectfont\color{shInkSecondary}
\textbf{Reinforcement Learning for Robotics --- The FCP Way}\par
\vspace{5pt}
Print edition, generated from the interactive web edition.\par
\vspace{9pt}
This book is built on four works, cited throughout as baseline references:
Sutton \& Barto's \emph{Reinforcement Learning: An Introduction} (2nd ed.) for the
mathematical spine; Kober, Bagnell \& Peters' \emph{Reinforcement Learning in
Robotics: A Survey} (IJRR 2013) for the robotics bridge; Tang et al.'s
\emph{Deep Reinforcement Learning for Robotics: A Survey of Real-World Successes}
(2024) for the modern taxonomy and the L0--L5 maturity rubric; and Akinola's
Columbia lecture notes for pedagogical framing.\par
\vspace{9pt}
Every chapter of the web edition carries interactive simulations that run the
real algorithms. In this print edition each appears as a described figure with
its widget identifier, so the two editions can be read side by side.\par
\vspace{9pt}
{\color{shInkMuted}Set in Space Grotesk and a Palatino-class serif.
Composed with \LaTeX.}\par}
\vspace*{\fill}

% ----- Contents --------------------------------------------------------------
\cleardoublepage
\pagestyle{fancy}
\setcounter{tocdepth}{1}
\renewcommand{\contentsname}{Contents}
\tableofcontents

% =============================================================================
%  Main matter
% =============================================================================
\mainmatter

\part{Foundations of Sequential Decision-Making}
\input{chapters/ch01}
\input{chapters/ch02}
\input{chapters/ch03}
\input{chapters/ch04}
\input{chapters/ch05}
\input{chapters/ch06}
\input{chapters/ch07}

\part{Scaling Up: Function Approximation and Deep RL}
\input{chapters/ch08}
\input{chapters/ch09}
\input{chapters/ch10}
\input{chapters/ch11}
\input{chapters/ch12}

\part{The Robotics Side}
\input{chapters/ch13}
\input{chapters/ch14}
\input{chapters/ch15}
\input{chapters/ch16}
\input{chapters/ch17}

\part{Competencies: RL on Real Robots}
\input{chapters/ch18}
\input{chapters/ch19}
\input{chapters/ch20}

\part{Frontiers and Capstone}
\input{chapters/ch21}
\input{chapters/ch22}

\end{document}
